\chapter{Ensemble Grandecanônico}

    O chamado \textit{Ensemble Grandecanônico} é definido considerando um sistema em contato com um reservatório térmico à temperatura $T$ e um reservatório de partículas com potencial químico $\mu$. Agora, olhamos para a \textit{distribuição} de estados $j$ segundo o \textit{peso de Boltzmann}, com a probabilidade de um estado sendo:
    \begin{equation}
        P(j) \propto e^{-E(j)\beta+N(j)\alpha}\,,
    \end{equation}
    com \[\beta=\frac{1}{k_BT}\,,\] e \[\alpha=\frac{\mu}{k_BT}\,.\] De fato, para conhecer a termodinâmica de um sistema no ensemble grandecanônico, basta conhecer a \textit{função de partição grandecanônica}

    \begin{statementbox}{Função de partição grandecanônica}
        Dado um ensemble grandecanônico, definimos sua função de partição $\mathcal{Z}_G$ como a soma dos pesos de Boltzmann sobre \textit{todos} os estados de equilíbrio possíveis ao sistema:
        \begin{equation}
            \mathcal{Z}_C(T,V,N) = \sum_{j}e^{-E(j)\beta+N(j)\alpha}.
        \end{equation}
        Esta função está ligada ao chamado \textit{potencial grande canônico} pela relação:
        \begin{equation}
            \mathcal{A}(T,V,\mu)=-k_BT\ln\mathcal{Z}_G.
        \end{equation}
    \end{statementbox}